\section{Local optimizer structure near the phase boundary}\label{sec:local-optimizer-structure}

We now state and prove the main local theorem in its precise form.
Recall the log-odds displacement
\[
  \lambda(p,r)=
  \log\frac{1-p}{p}-\log\frac{1-\pc(r)}{\pc(r)}.
\]
Thus, $\lambda(\pc(r),r)=0$, and $\lambda(p,r)>0$ is equivalent to
$p<\pc(r)$.

\begin{theorem}[Local structure of regular-graph optimizers]\label{thm:local-optimizer-structure}
Let $H$ be a fixed $d$-regular graph with $d\ge2$ and $m=|E(H)|$, and let
$r_0\in(0,1)\setminus\{r_*\}$.
Then, there exist $\rho>0$ and $\eta>0$ such that, for every
$r$ with $|r-r_0|<\rho$ and every $p$ with $0<p<r$ and
$|p-\pc(r)|<\eta$, the following hold.

\begin{enumerate}[label=(\alph*)]
  \item If $p\ge\pc(r)$, then the unique optimizer of
  \eqref{eq:graphon-variational-problem} is the constant graphon $W\equiv r$.
  \item If $p<\pc(r)$, then the optimizer is unique up to relabelling,
  nonconstant, and bipodal.
  \item If $p<\pc(r)$ and $W_{p,r}$ denotes the unique bipodal optimizer from
  (b), then
  \begin{equation}\label{eq:optimizer-edge-expansion}
    e(W_{p,r})=r-\frac{\lambda(p,r)}{A_H(r)}+O(\lambda(p,r)^2),
  \end{equation}
  and
  \begin{equation}\label{eq:optimal-value-expansion}
    \Phi_H(p,r)=J_p(r)-\frac{\lambda(p,r)^2}{2A_H(r)}+O(\lambda(p,r)^3),
  \end{equation}
  where the implicit constants in the \(O\)-terms are uniform over the above
  range \(|r-r_0|<\rho\), and \(A_H(r)\) is the positive analytic function
  from \Cref{thm:positive-second-variation} and \eqref{eq:second-variation-domain}.
  \item In the situation of \textup{(b)}, relabel $W_{p,r}$ so that its
  smaller pode is $A_{p,r}=[0,c(p,r)]$, where $0<c(p,r)<1/2$.  Then, there are
  densities $q_{11}(p,r)$, $q_{12}(p,r)$, and $q_{22}(p,r)$ such that
  \[
    W_{p,r}(x,y)=
    \begin{cases}
      q_{11}(p,r),
        & x,y\in A_{p,r},\\
      q_{12}(p,r),
        & (x,y)\in
          (A_{p,r}\times A_{p,r}^c)
          \cup(A_{p,r}^c\times A_{p,r}),\\
      q_{22}(p,r),
        & x,y\in A_{p,r}^c.
    \end{cases}
  \]
  The six functions
  \[
    \Phi_H(p,r),\quad e(W_{p,r}),\quad c(p,r),\quad
    q_{11}(p,r),\quad q_{12}(p,r),\quad q_{22}(p,r)
  \]
  are real-analytic in $(p,r)$ throughout the symmetry-breaking parameter
  region covered by the theorem.  Moreover, for every such $r$,
  \[
    c(p,r)\longrightarrow0
    \qquad\text{as }p\uparrow\pc(r).
  \]
\end{enumerate}
\end{theorem}

\begin{proof}
Choose a compact interval \(K\subset(0,1)\setminus\{r_*\}\) such that
\(r_0\in\operatorname{int}K\), and choose \(\rho>0\) so that
\((r_0-\rho,r_0+\rho)\subseteq K\).

First consider $p\ge\pc(r)$.
Condition (M2) of \Cref{thm:scalar-lz-boundary} is an equivalence valid for
\emph{every} $p\in(0,1)$, so it applies directly here and gives that
$(r^d,J_p(r))$ lies on the convex minorant of $x\mapsto J_p(x^{1/d})$.  Hence
the uniqueness clause of \Cref{thm:lz-criterion} gives that the
constant graphon $W\equiv r$ is the unique optimizer.  Note that this half of
the theorem needs no restriction on $p-\pc(r)$ at all: it holds for every
$p\in[\pc(r),r)$.

Now let $p<\pc(r)$.
Let $\rho_0$ and the function $\eta(\cdot)$ be as in
\Cref{cor:scalar-reduction}, and let $\overline\delta$, $a_0$, $C_2$
and the analytic extension $G(r,\delta)$ be as in
\Cref{thm:positive-second-variation}.  By \eqref{eq:boundary-excess-extension}, this function agrees
with the boundary excess \eqref{eq:boundary-excess} for
$0<\delta<\overline\delta$; at $\delta=0$, for negative $\delta$, and whenever
a $\delta$-derivative appears, we use the analytic extension $G$ itself.
Thus, \Cref{thm:positive-second-variation}(i) and (iii) read
\[
  G(r,0)=\partial_\delta G(r,0)=0,
  \qquad
  \partial_\delta^2G(r,\delta)=A_H(r)+O_K(|\delta|)
  \qquad
  (r\in K,\ |\delta|<\overline\delta),
\]
the implied constant being $C_2$.

Both $\rho_0$ and $\overline\delta$ are produced by applying \Cref{thm:krrs-analytic-extension} on
$K$; replacing each by the smaller of the two, we may and do assume that a
single Kenyon--Radin--Ren--Sadun neighbourhood serves both.  Choose
\[
  0<\delta_0<
  \min\left\{\rho_0,\overline\delta,\frac{a_0}{2C_2},\min K\right\}.
\]
Then, uniformly for \(r\in K\) and \(|\delta|<\delta_0\),
\[
  \partial_\delta^2G(r,\delta)\ge A_H(r)-C_2|\delta|
  \ge \frac{a_0}{2}>0.
\]
Choose $0<\eta<\eta(\delta_0)$.

For any symmetry-breaking pair $(p,r)$ with $r\in K$, $p<\pc(r)$, and
$|p-\pc(r)|<\eta$, \Cref{cor:scalar-reduction}, applied with
$\rho=\delta_0$ and $\eps=r-\delta$, reduces the variational problem defining
$\Phi_H(p,r)$ to
\[
  \Phi_H(p,r)
  =
  \min_{0<\delta<\delta_0}
  \mathcal F_{p,r}(r-\delta).
\]
Moreover, \Cref{cor:scalar-reduction} identifies the scalar
minimizers with the optimizers of the original graphon problem.  Every graphon
optimizer \(W^*\) is, up to relabelling, of the form
\[
  W^*=\widehat W_{r-\delta,r}
\]
for some minimizer \(\delta\in(0,\delta_0)\) of
\(\delta\mapsto\mathcal F_{p,r}(r-\delta)\), namely
\(\delta=r-e(W^*)\).  Conversely, for every such scalar minimizer \(\delta\),
the graphon \(\widehat W_{r-\delta,r}\) is an optimizer of the original
problem.  Since distinct values of \(\delta\) give distinct edge densities,
the graphon optimizer is unique up to relabelling if and only if the scalar
problem has a unique minimizer.
At the boundary value $p_0=\pc(r)$, we have, by \eqref{eq:boundary-excess} together
with \eqref{eq:boundary-excess-extension},
\[
  G(r,\delta)
  =
  \mathcal F_{p_0,r}(r-\delta)-J_{p_0}(r)
  \qquad(0<\delta<\delta_0).
\]
We now make explicit how changing \(p\) to \(p_0\) affects the reduced
objective.  The relative entropy identity
\eqref{eq:relative-entropy-identity} gives, for any graphon \(W\),
\[
  I_p(W)-I_{p_0}(W)
  =
  \log\frac{1-p_0}{1-p}
  +e(W)\left(
    \log\frac{1-p}{p}-\log\frac{1-p_0}{p_0}
  \right).
\]
The expression in parentheses is \(\lambda(p,r)\).  Since
\(\mathcal F_{p,r}(r-\delta)\) is the \(I_p\)-value of the fixed-$(e,t(H,\cdot))$
entropy maximizer \(\widehat W_{r-\delta,r}\), whose edge density is \(r-\delta\),
we obtain
\[
  \mathcal F_{p,r}(r-\delta)
  -\mathcal F_{p_0,r}(r-\delta)
  =
  \log\frac{1-p_0}{1-p}
  +(r-\delta)\lambda(p,r).
\]
Applying the same identity to the constant graphon \(W\equiv r\) gives
\[
  J_p(r)-J_{p_0}(r)
  =
  \log\frac{1-p_0}{1-p}
  +r\lambda(p,r).
\]
Subtracting the last two identities cancels the logarithmic term and the
\(r\lambda(p,r)\) term, leaving
\[
  \bigl(\mathcal F_{p,r}(r-\delta)-J_p(r)\bigr)
  -\bigl(\mathcal F_{p_0,r}(r-\delta)-J_{p_0}(r)\bigr)
  =-\lambda(p,r)\delta.
\]
Consequently,
\[
  \mathcal F_{p,r}(r-\delta)-J_p(r)
  =
  \bigl(\mathcal F_{p_0,r}(r-\delta)-J_{p_0}(r)\bigr)
  -\lambda(p,r)\delta
  =
  G(r,\delta)-\lambda(p,r)\delta.
\]
Therefore, \eqref{eq:graphon-variational-problem} reduces to the scalar problem
\begin{equation}\label{eq:edge-deficit-minimization}
  \Phi_H(p,r)
  =\min_{0<\delta<\delta_0}
  \{G(r,\delta)-\lambda(p,r)\delta\}
  + J_p(r).
\end{equation}
The minimum over the open interval is attained.  To see this, the graphon
variational problem \eqref{eq:graphon-variational-problem} has an optimizer: the
constant graphon $W\equiv r$ shows that its feasible set is nonempty, and
graphon compactness together with the lower semicontinuity of $I_p$ gives
attainment.  By \Cref{cor:scalar-reduction}, any such optimizer
$W^*$ yields a scalar minimizer
$\delta=r-e(W^*)\in(0,\delta_0)$.

We now locate the critical point of the scalar objective
\(\delta\mapsto G(r,\delta)-\lambda(p,r)\delta\) and show that it varies
analytically with \((p,r)\).  Its critical-point equation is
\[
  \Psi(p,r,\delta)
  :=
  \partial_\delta G(r,\delta)-\lambda(p,r)
  =0.
\]
The joint analyticity of \(G(r,\delta)\), proved in
\Cref{thm:positive-second-variation}, together with the analyticity of
\(\lambda(p,r)\), implies that \(\Psi\) is analytic in a neighbourhood of the
lifted phase-boundary curve
\[
  \{(\pc(r),r,0):r\in K\}.
\]
For every \(\bar r\in K\),
\[
  \Psi(\pc(\bar r),\bar r,0)
  =
  \partial_\delta G(\bar r,0)-\lambda(\pc(\bar r),\bar r)
  =0,
  \qquad
  \partial_\delta\Psi(\pc(\bar r),\bar r,0)
  =A_H(\bar r)>0.
\]
The analytic implicit function theorem, with \((p,r)\) as parameters and
\(\delta\) as the unknown, therefore gives a neighbourhood
\(U_{\bar r}\) of \((\pc(\bar r),\bar r)\) and a unique analytic function
\(\delta_{\bar r}:U_{\bar r}\to(-\delta_0,\delta_0)\) satisfying
\[
  \Psi(p,r,\delta_{\bar r}(p,r))=0,
  \qquad
  \delta_{\bar r}(\pc(\bar r),\bar r)=0.
\]
In other words, the critical point \(\delta=0\), which the scalar objective has
at the phase-boundary point \((\pc(\bar r),\bar r)\), persists as \((p,r)\)
moves away from that point and depends analytically on \((p,r)\).

These local functions patch together over the parameter range of interest.
To see this, for fixed \((p,r)\) with \(r\in K\),
\[
  \partial_\delta\Psi(p,r,\delta)=\partial_\delta^2G(r,\delta)
  \ge\frac{a_0}{2}
  \qquad (|\delta|<\delta_0),
\]
so \(\Psi(p,r,\cdot)\) has at most one zero in
\((-\delta_0,\delta_0)\).  Hence, two local functions agree on the relevant
part of any overlap.  The compact phase-boundary curve
\[
  \{(\pc(r),r):r\in K\}
\]
is covered by finitely many of the \(U_{\bar r}\).  After decreasing \(\eta\)
if necessary, their union contains every \((p,r)\) with \(r\in K\) and
\(|p-\pc(r)|<\eta\).  The patched functions therefore define a single
analytic function \(\delta_*(p,r)\) on this uniform neighbourhood, with
  \[
  \delta_*(\pc(r),r)=0,
  \qquad
  \partial_\delta G(r,\delta_*(p,r))=\lambda(p,r).
  \]

Since $p<\pc(r)$, we have $\lambda(p,r)>0$; the critical-point equation and
$\partial_\delta G(r,0)=0$ therefore imply that $\delta_*(p,r)\ne0$.
The mean value theorem now gives
\[
  \lambda(p,r)
  =
  \partial_\delta^2G(r,\theta\delta_*(p,r))\delta_*(p,r)
\]
for some \(\theta\in(0,1)\).  Since
\(\partial_\delta^2G(r,\cdot)\ge a_0/2\) on the relevant interval,
\(\delta_*(p,r)\) has the same sign as \(\lambda(p,r)\) and
\[
  |\delta_*(p,r)|
  \le \frac{2}{a_0}|\lambda(p,r)|.
\]
Moreover,
\(\partial_\delta^2G(r,\theta\delta_*)=A_H(r)+O_K(|\delta_*|)\).  Hence, the preceding
identity and bound give
\[
  \delta_*-\frac{\lambda(p,r)}{A_H(r)}
  =\frac{\lambda(p,r)}{\partial_\delta^2G(r,\theta\delta_*)}
  -\frac{\lambda(p,r)}{A_H(r)}
  =
  \frac{\lambda(p,r)
  \bigl(A_H(r)-\partial_\delta^2G(r,\theta\delta_*)\bigr)}
  {\partial_\delta^2G(r,\theta\delta_*)A_H(r)}
  =
  \frac{O_K(\lambda(p,r)^2)}
  {\partial_\delta^2G(r,\theta\delta_*)A_H(r)}
  =O_K(\lambda(p,r)^2).
\]
Here the denominator is at least \(a_0^2/2\).  Thus,
\begin{equation}\label{eq:optimal-deficit-expansion}
  \delta_*(p,r)
  =
  \frac{\lambda(p,r)}{A_H(r)}+O_K(\lambda(p,r)^2).
\end{equation}
The sign conclusion above gives $\delta_*(p,r)>0$.  By construction,
$\delta_*(p,r)\in(-\delta_0,\delta_0)$, and hence
\[
  0<\delta_*(p,r)<\delta_0.
\]
Thus, $\delta_*(p,r)$ lies in the range over which the minimum in
\eqref{eq:edge-deficit-minimization} is taken.
Since $\partial_\delta^2G(r,\delta)>0$ on the whole interval
$(0,\delta_0)$, the scalar objective
$G(r,\delta)-\lambda(p,r)\delta$ is strictly convex on that whole interval.
So $\delta_*(p,r)$ is the unique minimizer of
$\delta\mapsto G(r,\delta)-\lambda(p,r)\delta$ on $(0,\delta_0)$, and the
minimum in \eqref{eq:edge-deficit-minimization} is attained there.

By the preceding reduction, every optimizer attaining $\Phi_H(p,r)$ in
\eqref{eq:graphon-variational-problem} is, up to relabelling, the fixed-$(e,t(H,\cdot))$
entropy maximizer $\widehat W_{r-\delta_*,r}$ supplied by
\Cref{thm:krrs-analytic-extension}, with
\[
  e(\widehat W_{r-\delta_*,r})=r-\delta_*,
  \qquad
  t(H,\widehat W_{r-\delta_*,r})=r^m.
\]
Writing this optimizer as $W_{p,r}$, the graphon $W_{p,r}$ is the unique optimizer up to
relabelling and is bipodal.  It is nonconstant because it has edge density
\(r-\delta_*<r\) while its \(H\)-density is \(r^m\), whereas a constant graphon
with \(H\)-density \(r^m\) has edge density \(r\).
Write $\widetilde c(\eps,\vartheta)$ and
$\widetilde q_{ij}(\eps,\vartheta)$ for the analytic KRR--S parameter maps in
\Cref{thm:krrs-analytic-extension}, and set
\[
  \eps(p,r):=r-\delta_*(p,r),
  \qquad
  \vartheta(p,r):=r^m-\eps(p,r)^m.
\]
Thus, the parameters in \textup{(d)} are defined by
\[
  c(p,r):=\widetilde c\bigl(\eps(p,r),\vartheta(p,r)\bigr),
  \qquad
  q_{ij}(p,r):=
  \widetilde q_{ij}\bigl(\eps(p,r),\vartheta(p,r)\bigr)
  \quad(ij\in\{11,12,22\}).
\]
These four functions are analytic in $(p,r)$.  The identity
$e(W_{p,r})=r-\delta_*(p,r)$ proves the same for the edge density, and
\eqref{eq:edge-deficit-minimization}, evaluated at
$\delta=\delta_*(p,r)$, proves it for $\Phi_H(p,r)$.

Moreover, $\lambda$ is continuous and vanishes on the compact phase-boundary
curve $\{(\pc(r),r):r\in K\}$.  Thus, $\lambda(p,r)\to0$ as
$p\uparrow\pc(r)$, uniformly for $r\in K$, and the preceding bound on
$\delta_*$ gives
\[
  \delta_*(p,r)\longrightarrow0
  \qquad\text{as }p\uparrow\pc(r),
\]
uniformly for $r\in K$.  Hence,
$(\eps(p,r),\vartheta(p,r))\to(r,0)$ uniformly.  Since
$\widetilde c(r,0)=0$ by \Cref{thm:krrs-analytic-extension} and the extended parameter map is
continuous near the compact set $K\times\{0\}$, it follows that
$c(p,r)\to0$ uniformly for $r\in K$.  Since $\delta_*(p,r)>0$, we have
$\eps(p,r)<r$ and hence $\vartheta(p,r)>0$.  For this positive surplus, the
normalization in \Cref{thm:krrs-analytic-extension} makes the first pode the smaller one; that
pode has positive measure, since a null first pode would give $H$-density
$\eps(p,r)^m$ rather than $r^m$.  Hence, $c(p,r)>0$; after decreasing $\eta$
if necessary, the
uniform convergence also gives $c(p,r)<1/2$ throughout the symmetry-breaking
region.  A
measure-preserving relabelling identifies this first pode with
$A_{p,r}=[0,c(p,r)]$; the three parameter functions then give the values of
$W_{p,r}$ on $A_{p,r}\times A_{p,r}$, on the two cross rectangles, and on
$A_{p,r}^c\times A_{p,r}^c$, respectively, as stated in \textup{(d)}.
In particular, the edge density has the expansion
\[
  e(W_{p,r})=r-\delta_*(p,r)
  =r-\frac{\lambda(p,r)}{A_H(r)}+O_K(\lambda(p,r)^2),
\]
and, since $\delta_*>0$, \eqref{eq:boundary-excess-extension} and
\eqref{eq:boundary-excess-expansion} give
\[
  G(r,\delta_*)=\frac12A_H(r)\delta_*^2+O_K(\delta_*^3)
\]
together with \eqref{eq:optimal-deficit-expansion}, the objective value has the expansion
\begin{align*}
  \Phi_H(p,r)
  =\left(J_p(r)+G(r,\delta_*)-\lambda(p,r)\delta_*\right)
  =\left(J_p(r)-\frac{\lambda(p,r)^2}{2A_H(r)}+O_K(\lambda(p,r)^3)\right).
\end{align*}
For the last step, write $\delta_*=\lambda(p,r)/A_H(r)+e$ with
$e=O_K(\lambda(p,r)^2)$ by \eqref{eq:optimal-deficit-expansion}; then
\[
  \frac12A_H(r)\delta_*^2-\lambda(p,r)\delta_*
  =-\frac{\lambda(p,r)^2}{2A_H(r)}+\frac12A_H(r)e^2,
\]
and $\frac12A_H(r)e^2=O_K(\lambda(p,r)^4)=O_K(\lambda(p,r)^3)$ because
$\lambda(p,r)$ is bounded on the range under consideration, while
$O_K(\delta_*^3)=O_K(\lambda(p,r)^3)$ by $\delta_*\le(2/a_0)\lambda(p,r)$.
\end{proof}

First-order refinements for the individual bipodal parameters are derived in
Appendix~\ref{app:bipodal-parameter-expansions}.